\label{sec:framework}

\subsection{Character Vectors}

Let $\mathcal{T}$ denote the set of census tracts in the state under redistricting.
For each tract $t \in \mathcal{T}$, we define a \emph{community character vector}
\begin{equation}
  \mathbf{c}(t) = \bigl(c_1(t),\, c_2(t),\, \ldots,\, c_k(t)\bigr) \in \mathbb{R}^k_{\geq 0},
  \label{eq:char-vector}
\end{equation}
where each coordinate $c_i(t)$ is a non-negative real-valued signal derived from
publicly available, partisan-neutral data.
The number of dimensions $k$ depends on which signals are active for the given plan
configuration; any subset of M.1--M.7 may contribute one or more dimensions.
We require only that each $c_i(t) \geq 0$ and that the vector is defined for every
tract with nonzero population.

\subsection{Cosine Similarity}

Given two adjacent tracts $u, v \in \mathcal{T}$, their community similarity is
\begin{equation}
  \mathrm{sim}(\mathbf{c}_u, \mathbf{c}_v)
  \;=\;
  \frac{\mathbf{c}_u \cdot \mathbf{c}_v}{\|\mathbf{c}_u\| \, \|\mathbf{c}_v\|}
  \;\in\; [0,\, 1],
  \label{eq:cosine-sim}
\end{equation}
where the dot product and norms are the standard Euclidean ones.
Because all character vector coordinates are non-negative, the cosine similarity is
always in $[0,1]$; negative similarity---which would arise from negative character
signals and would have no natural interpretation---is impossible by construction.

Cosine similarity measures the \emph{orientation} of the character vectors rather
than their magnitude.
This choice is deliberate: two tracts that are economically active in the same mix of
industries should be recognized as similar whether one is a large urban tract and
the other is a small suburban one.
Scaling by population or area is not done at the similarity level; if population
weighting is desired for a particular signal, it enters $c_i(t)$ before the similarity
computation, not after it.

\subsection{Edge Weight Modifier}

Let $G = (\mathcal{T}, E)$ be the tract adjacency graph, where an edge
$(u,v) \in E$ exists when tracts $u$ and $v$ share a boundary of positive length
(zero-length point adjacencies are excluded).
The \textsc{bisect} platform assigns each edge a base weight $w_{\mathrm{base}}(u,v)$
determined by the active base weight mode (geographic boundary length, county
co-membership, or unweighted).
The community character modifier replaces or augments that base weight as
\begin{equation}
  w(u,v) \;=\; w_{\mathrm{base}}(u,v) \;\times\; \mathrm{sim}(\mathbf{c}_u, \mathbf{c}_v).
  \label{eq:weight-modifier}
\end{equation}
When $\mathrm{sim}(\mathbf{c}_u, \mathbf{c}_v) = 1$ (identical character vectors),
the weight is unchanged from $w_{\mathrm{base}}$.
When $\mathrm{sim}(\mathbf{c}_u, \mathbf{c}_v) = 0$ (orthogonal character vectors,
meaning completely different community profiles), the edge weight collapses to zero,
making the cut between those tracts free.
Intermediate similarities yield intermediate weights.

\subsection{Mathematical Properties}

\paragraph{Symmetry.}
$\mathrm{sim}(\mathbf{c}_u, \mathbf{c}_v) = \mathrm{sim}(\mathbf{c}_v, \mathbf{c}_u)$
for all pairs $(u,v)$, since the dot product is symmetric.
Consequently $w(u,v) = w(v,u)$, which METIS requires: undirected graphs have
symmetric edge weights.

\paragraph{Non-negativity.}
$\mathrm{sim}(\mathbf{c}_u, \mathbf{c}_v) \geq 0$ whenever
$\mathbf{c}_u, \mathbf{c}_v \in \mathbb{R}^k_{\geq 0}$, since the dot product of
two non-negative vectors is non-negative.
METIS requires non-negative integer edge weights; the implementation rounds
$w(u,v)$ to the nearest positive integer after scaling, with a minimum value of $1$.

\paragraph{Graceful degradation.}
If all tracts have identical character vectors---formally, if
$\mathbf{c}(t) = \mathbf{c}$ for all $t$---then
$\mathrm{sim}(\mathbf{c}_u, \mathbf{c}_v) = 1$ for all edges, and
$w(u,v) = w_{\mathrm{base}}(u,v)$ for all edges.
The community character modifier does not change the plan at all when there is no
character variation to exploit.
This ensures that the framework introduces no distortion in states where the chosen
signal has no geographic variation.

\paragraph{Degenerate case: zero vector.}
If a tract $t$ has $\mathbf{c}(t) = \mathbf{0}$---for example, because no workers
reside in a tract that is entirely commercial---the cosine similarity is undefined.
The implementation handles this by substituting $\mathrm{sim} = 1$ (treat as identical
to any neighbor), which is conservative: it does not create a free cut but also does
not create an artificially strong bond.

\subsection{Well-Definedness of the Similarity}

Cosine similarity is undefined when either vector is the zero vector.
A tract with zero character in all dimensions is a tract that the available data
cannot characterize; it should not bias the partitioner in either direction.
The substitution $\mathrm{sim} = 1$ for zero-vector tracts is equivalent to treating
such tracts as perfectly similar to their neighbors in the character dimension,
leaving their cut costs governed entirely by the base weight mode.

Additionally, for the administrative zone co-membership signal (M.6), similarity is
not cosine-based but ratio-based (fraction of zones shared); see M.6 for the precise
formula.
The zone co-membership formula satisfies symmetry, $[0,1]$ range, and graceful
degradation independently, and is a valid plug-in replacement for
$\mathrm{sim}(\mathbf{c}_u, \mathbf{c}_v)$ in equation~(\ref{eq:weight-modifier}).

\subsection{Invariants for Implementation}

Any correct implementation of the framework must satisfy:
\begin{itemize}
  \item \texttt{similarity\_symmetric}: $\mathrm{sim}(u,v) = \mathrm{sim}(v,u)$
    for all adjacent tract pairs.
  \item \texttt{similarity\_self\_is\_one}: $\mathrm{sim}(u,u) = 1.0$ for all
    tracts $u$ (a tract is perfectly similar to itself).
  \item \texttt{weight\_stays\_nonneg}: $w(u,v) \geq 0$ for all edges after
    blending and rounding.
  \item \texttt{weight\_bounded\_above}: $w(u,v) \leq 2 \times w_{\mathrm{base}}(u,v)$
    under the multiplicative formula (since $\mathrm{sim} \leq 1$).
\end{itemize}

These are L0 unit test invariants checked on every build of the \textsc{bisect}
crate implementing this framework.
